\section{Materials and Methods}
\label{sec:materials_and_methods}

\begin{figure*}[htbp]
\centering
\includegraphics[width=\textwidth]{method.png}
\caption{End-to-End Framework of the Proposed Multi-Domain Vision Transformer Feature Fusion and Classical Classifier Optimization for Intersectional Race and Gender Classification.}
\label{fig:framework}
\end{figure*}

Kerangka kerja yang diusulkan memadukan representasi visual multi-domain dan optimasi model pembelajaran mesin untuk klasifikasi ras dan gender interseksional, sebagaimana diilustrasikan pada Figure~\ref{fig:framework}. Metodologi diawali pemisahan citra wajah $224 \times 224$ piksel pada dataset DemogPairs menjadi partisi latih dan uji via stratified split. Tiga model ViT-Base pre-trained yang dibekukan mengekstraksi representasi laten berdimensi 768 dari token [CLS] secara offline. Ketiga vektor digabungkan melalui fusi konkatenasi $\mathbf{z}_{\text{tri}} = \mathbf{f}_{\text{face}} \oplus \mathbf{f}_{\text{emotion}} \oplus \mathbf{f}_{\text{age}}$, di mana $\mathbf{z}_{\text{tri}}$ menyatakan vektor tri-domain berdimensi 2,304, $\oplus$ adalah operator konkatenasi, serta $\mathbf{f}_{\text{face}}$, $\mathbf{f}_{\text{emotion}}$, dan $\mathbf{f}_{\text{age}}$ merepresentasikan fitur biometrik wajah, ekspresi, dan usia. Vektor representasi diproses dalam pipeline modular Scaler dan PCA sebelum diklasifikasikan ke enam kelas interseksional oleh RF, GNB, LR, dan SVM. Penalaan hyperparameter dijalankan melalui 5-Fold Stratified Cross-Validation dengan skema GridSearchCV terisolasi untuk mencegah kebocoran informasi sebelum evaluasi data uji.
